\section{Carthage and the Making of a Latin Christian Voice}

\subsection{A metropolis, not a provincial margin}

The Carthage from which Tertullian writes was the rebuilt Roman capital of
Africa Proconsularis: a Mediterranean port, administrative center, arena of
imperial display, and home to schools of grammar and rhetoric. Its inhabitants
moved among Latin, Greek, and Punic worlds; wealth, enslavement, patronage,
trade, municipal ambition, and older regional identities occupied the same
city. To call Tertullian ``African'' securely locates him in this province and
its Christian networks. It does not by itself establish Punic or Berber
ancestry, an anti-Roman national consciousness, or a single African theology.
Those are modern historical proposals that must be argued from more than
geography.

Christianity was present in Latin-speaking North Africa before his surviving
career. The acts of the Scillitan martyrs record a proconsular hearing at
Carthage on 17 July 180 in which Christians from Scillium refused sacrifice and
were executed. The record is brief and stylized in transmission, but it anchors
a Latin Christian community and its exposure to Roman justice. During
Tertullian's literary activity, the prison diary and martyr narrative associated
with Perpetua, Felicity, and their companions place catechumens, visions,
baptism, prison ministry, conflict within a family, and execution at Carthage
in 202 or 203. The old suggestion that Tertullian edited the \emph{Passion of
Perpetua and Felicity} is possible literary history, not an established event
in his life.

The result is a more demanding setting than the phrase ``first Latin
theologian'' sometimes suggests. Tertullian did not invent Latin Christianity
alone. He inherited communities already confessing, reading, baptizing,
disciplining, and dying in Latin. What survives from him is the first large
Latin Christian authorial corpus, not the first Christian thought ever formed
in the language.

\subsection{The adult author appears}

No secure childhood scene precedes the books. Tertullian displays formidable
training in rhetoric, argument, literary parody, Roman antiquarian material,
philosophical schools, medicine, and legal idiom. The inference that he received
an elite literary education is strong because it rests on repeated performance,
not on Jerome's later biography. It does not prove that he held a professional
appointment as advocate or that he was identical with a jurist named
Tertullianus cited in the \emph{Digest}. Legal competence, a Roman law career,
and residence at Rome are three different claims.

Eusebius, writing at Caesarea more than a century later, calls him especially
versed in Roman laws and famous in other respects (\emph{Ecclesiastical
History} 2.2.4). Eusebius also knew a Greek translation of the
\emph{Apologeticum}, from which he quotes. That testimony proves the early
Greek circulation and reputation of the African author better than it proves a
lost curriculum vitae. Jerome's \emph{On Illustrious Men} 53, written about
392 or 393, supplies the famous claim that Tertullian was the son of a
\emph{centurio proconsularis}. No earlier located witness names that father,
and the phrase itself has generated competing attempts to identify an office.
It belongs to received biography, not the secure childhood.

Only one part of the expansive name often printed on book covers is anchored
inside an authorial work. The closing sentence of \emph{On the Veiling of
Virgins} 17 names ``Septimius Tertullianus'' in the third person and says the
little work is his. This transmitted closing attribution is the earliest
internal name form, but its syntax is not a modern signature. The fuller
``Quintus Septimius Florens Tertullianus'' comes through medieval paratext and
manuscript titles whose forms vary. Those witnesses matter for the reception of
his books; they do not all become a birth record.

\subsection{Writing from a political wound}

Two works allow an unusual chronological foothold. \emph{To the Nations}
1.17.4 alludes to the defeat of Clodius Albinus, whose forces lost to Septimius
Severus near Lugdunum on 19 February 197. The \emph{Apologeticum} reworks much
of the same case and is ordinarily placed later that year. These are not first
drafts of a private conversion diary. They are public literary interventions
by an already formed Christian writer.

Carthage had reason to feel the struggle. Albinus had drawn support in the
western provinces, and imperial victory brought punishments and political
realignment. Tertullian exploits memories of governors, senatorial conspiracy,
imperial succession, crowd accusation, and provincial trial. The context helps
explain the temperature of the prose, but it does not reduce the
\emph{Apologeticum} to partisan propaganda or prove an empire-wide program of
continuous ``Severan persecution.'' Christians were vulnerable through local
accusation, judicial choice, public ritual, and intermittent coercion. The
archive requires local events and legal rhetoric, not one smooth persecution
timeline.

Already the durable pattern is visible. Tertullian asks institutions to become
what their stated rules imply; he exposes a gap between principle and practice;
and he answers ambiguity by sharpening the boundary. This made him an
extraordinary advocate for Christians. It also supplied the habit by which he
would later prosecute Christians who accepted disciplines he judged too easy.
